\chapter*{Lời mở đầu}
\addcontentsline{toc}{chapter}{Lời mở đầu}

\songngu{Bạn nói rất đúng: bốn cuốn mỏng, mỗi cuốn một chút, đôi khi đọc bị đứt mạch.}{You were right: four thin books, each giving only a little, can feel broken in flow.}
\songngu{Vì vậy cuốn này gộp luôn Quyển 38, 39, 40 và 41 vào một quyển duy nhất, nhưng không trộn lẫn lung tung mà chia thành đúng 4 phần.}{So this book merges Volumes 38, 39, 40, and 41 into one single volume, but instead of mixing everything together, it is divided into four exact parts.}
\songngu{Điểm mới quan trọng hơn là toàn bộ nội dung được viết theo kiểu song ngữ từng câu.}{The more important upgrade is that the whole book is written in sentence-by-sentence bilingual form.}
\songngu{Bạn có thể đọc câu tiếng Việt trước, nhìn ngay câu tiếng Anh bên dưới, rồi nhặt từng cụm để học luôn.}{You can read the Vietnamese sentence first, look at the English sentence right below it, and pick up phrases as you go.}

\songngu{Phần I giữ tinh thần của Quyển 38: cứu tiết học đang trôi và kéo nhịp lớp quay lại.}{Part I keeps the spirit of Volume 38: rescuing a drifting lesson and pulling the class rhythm back.}
\songngu{Phần II đi theo Quyển 39: đối đáp có duyên, giữ không khí mà vẫn giữ lớp.}{Part II follows Volume 39: witty dialogue that keeps the atmosphere while still keeping the class in hand.}
\songngu{Phần III lấy trục của Quyển 40: hỏi để kéo cả lớp nói và nghĩ.}{Part III takes the axis of Volume 40: asking questions that pull the whole class into speaking and thinking.}
\songngu{Phần IV đi theo Quyển 41: khởi động lớp thật nhanh trong vài phút đầu hoặc sau một cú tụt năng lượng.}{Part IV follows Volume 41: warming up the class quickly in the first few minutes or after an energy drop.}

\songngu{Cuốn này không cố làm học thuật.}{This book does not try to be academic.}
\songngu{Nó cố làm ba việc thôi: vui hơn, dùng được ngay hơn, và giúp bạn có thêm tiếng Anh thực dụng trong ngữ cảnh lớp học.}{It tries to do only three things: be more fun, be more usable right away, and give you practical English inside classroom situations.}
\songngu{Nếu một câu làm lớp cười nhẹ, vào bài nhanh, và đồng thời cho bạn nhớ thêm một cụm tiếng Anh, thế là cuốn sách đã làm đúng việc.}{If a line makes the class smile a little, enter the lesson faster, and helps you remember one more English phrase, then the book has done its job.}